\subsection{The Moment of Truth}
\label{subsec:11-2-moment-of-truth}

%%%%%%%%%%%%%%%%%%%%%%%%%%%%%%%%%%%%%%%%%%%%%%%%%%%%%%%%%%%%%%%%%%%%%%%%%%%%%%%
% SECTION 11.2: The Moment of Truth
% Module: BCM2_05_AWX_7240
% Content: t*, Situativity, Contextual Activation
%%%%%%%%%%%%%%%%%%%%%%%%%%%%%%%%%%%%%%%%%%%%%%%%%%%%%%%%%%%%%%%%%%%%%%%%%%%%%%%

The Moment of Truth ($t^*$) is the central temporal concept of the Awareness
function. Awareness exists only at this moment---not before, not after.

%------------------------------------------------------------------------------
\subsubsection{Definition}
%------------------------------------------------------------------------------

\begin{definition}[Moment of Truth]
\label{def:moment-of-truth}
The \textbf{Moment of Truth} ($t^*$) is the decision point at which potential 
utility can be transformed into behavior through awareness.
\end{definition}

\paragraph{Formal Condition.}
\begin{equation}
A_X(t^*) > 0 \quad \Leftrightarrow \quad \text{Decision active} \wedge \text{Context salience} > \theta_{awareness}
\end{equation}

\paragraph{Implication.}
Awareness can only operate at $t^*$, not before or after. This distinguishes
awareness from knowledge, attitudes, or intentions, which can exist across time.

%------------------------------------------------------------------------------
\subsubsection{Property 1: Situativity}
%------------------------------------------------------------------------------

Awareness operates \emph{only} at $t^*$, not before or after:

\begin{equation}
A_X(t) = 0 \quad \text{for all } t \neq t^*
\end{equation}

This distinguishes awareness from related constructs:

\begin{table}[htbp]
\centering
\caption{Awareness vs. Related Constructs}
\label{tab:awareness-vs-constructs}
\renewcommand{\arraystretch}{1.3}
\begin{tabular}{@{}llll@{}}
\toprule
\textbf{Construct} & \textbf{Temporal Scope} & \textbf{Situational?} & \textbf{Behavior-Guiding?} \\
\midrule
Knowledge & Persistent & No & Only if activated \\
Attitudes & Stable & No & Indirect \\
Intentions & Future-oriented & Partial & Indirect \\
\textbf{Awareness} & \textbf{Only at $t^*$} & \textbf{Yes} & \textbf{Direct} \\
\bottomrule
\end{tabular}
\end{table}

%------------------------------------------------------------------------------
\subsubsection{Property 2: Contextual Activation}
%------------------------------------------------------------------------------

Awareness at $t^*$ depends on the current context configuration:

\begin{equation}
A_{SYS_i}(t^*) = f(C(t^*), Norm(t^*), Emotion(t^*), \sigma_U(t^*))
\end{equation}

where:
\begin{itemize}
\item $C(t^*)$: Current context configuration (from $\Psi$-dimensions)
\item $Norm(t^*)$: Active social norms
\item $Emotion(t^*)$: Current emotional state
\item $\sigma_U(t^*)$: Utility salience
\end{itemize}

\paragraph{Context Determines Awareness.}
The same person with the same knowledge can have completely different
awareness depending on context:

\begin{table}[htbp]
\centering
\caption{Context-Dependent Awareness Activation}
\label{tab:context-awareness}
\renewcommand{\arraystretch}{1.3}
\begin{tabular}{@{}lll@{}}
\toprule
\textbf{Context} & \textbf{Typical $A$ Pattern} & \textbf{Example} \\
\midrule
Calm deliberation & High, balanced across categories & Financial planning session \\
Time pressure & High $A^{self}_{INU,t=0}$; low others & Rushed purchase decision \\
Social observation & High $A^{self}_{KNU}$, $A^{self}_{IDN}$ & Public behavior \\
Emotional arousal & Category-specific spikes & Fear → $A^{self}_{INU,P}$ \\
Cognitive depletion & Default to $A^{self}_{INU}$ & End of workday \\
\bottomrule
\end{tabular}
\end{table}

%------------------------------------------------------------------------------

%------------------------------------------------------------------------------
\subsubsection{Property 2a: Context Modulation (AWX-A7a)}
%------------------------------------------------------------------------------

Building on contextual activation, AWX-A7a specifies that context can either 
\emph{amplify} or \emph{suppress} awareness:

\begin{axiom}[AWX-A7a: Kontext-Modulation]
Context can amplify or suppress awareness at the Moment of Truth.

\textbf{Formal Specification:}
\begin{equation}
A_{mod} = f(C(t))
\end{equation}
where $f$ can take positive (amplification) or negative (suppression) values.
\end{axiom}

\paragraph{Amplification Examples.}
\begin{itemize}
    \item Health warning at point of purchase $\rightarrow$ raises $A_{INU}^{health}$
    \item Social presence of peers $\rightarrow$ amplifies $A_{IDN}^{reputation}$
    \item Deadline proximity $\rightarrow$ increases temporal awareness
\end{itemize}

\paragraph{Suppression Examples.}
\begin{itemize}
    \item Cognitive load from unrelated task $\rightarrow$ reduces all $A_X$
    \item Strong positive emotion $\rightarrow$ suppresses risk awareness
    \item Social conformity pressure $\rightarrow$ suppresses individual $A_{INU}$
\end{itemize}

\paragraph{Design Implication.}
Behavioral interventions should consider both directions. A poorly designed 
nudge may inadvertently suppress awareness of important considerations while 
amplifying less relevant ones. The Ψ-context dimensions (Section 11.16) 
systematically capture these modulation effects.

\subsubsection{Property 3: No Automatic Persistence}
%------------------------------------------------------------------------------

Awareness does not automatically persist across moments:

\begin{equation}
A_X(t^* + \Delta t) \neq A_X(t^*)
\end{equation}

Each new Moment of Truth requires fresh awareness activation. This explains:
\begin{itemize}
\item Why New Year's resolutions fade (awareness at midnight $\neq$ awareness in February)
\item Why reminders are needed (awareness must be re-activated)
\item Why habits bypass awareness (habitualized behavior requires no $A$)
\end{itemize}

%------------------------------------------------------------------------------
\subsubsection{The Moment of Truth in Different Contexts}
%------------------------------------------------------------------------------

\begin{table}[htbp]
\centering
\caption{Moment of Truth Across Domains}
\label{tab:moment-domains}
\renewcommand{\arraystretch}{1.3}
\begin{tabular}{@{}lll@{}}
\toprule
\textbf{Domain} & \textbf{Moment of Truth} & \textbf{Duration} \\
\midrule
Consumer purchase & At checkout/click & Seconds \\
Health behavior & When craving peaks & Seconds \\
Financial decision & When allocating funds & Minutes \\
Career transition & When accepting/declining offer & Hours \\
Policy vote & In voting booth & Minutes \\
Relationship decision & During conflict/proposal & Minutes to hours \\
\bottomrule
\end{tabular}
\end{table}

%------------------------------------------------------------------------------
\subsubsection{Axiom Reference: AWX-A2}
%------------------------------------------------------------------------------

\begin{axiom}[AWX-A2: Situativity]
\label{ax:awx-a2-ref}
Awareness operates exclusively at the moment of decision.
\end{axiom}

\begin{equation}
A_X(t) = 0 \quad \text{for all } t \neq t^*
\end{equation}

\paragraph{Sub-axioms:}
\begin{itemize}
\item \textbf{AWX-A2a (Situational Activation):} Only context-specific triggered 
awareness is effective.
\begin{equation}
A_{mod}(t) = 1 \text{ only when } Trigger_{ctx}(t) = true
\end{equation}

\item \textbf{AWX-A2b (General Awareness as Learning Effect):} General awareness 
arises exclusively through repeated situational activation.
\begin{equation}
A_{gen}(t) = \sum_{\tau=0}^{t} \delta A_{mod}(\tau)
\end{equation}

\item \textbf{AWX-A2c (Dynamic Salience):} Awareness changes depending on the 
temporal dynamics of perceived context.
\begin{equation}
\sigma_{dyn}(C_i, t) = \sigma_0(C_i) + \rho \cdot \frac{dC_i}{dt}
\end{equation}
\end{itemize}

%------------------------------------------------------------------------------
\subsubsection{Implications for Intervention}
%------------------------------------------------------------------------------

Because awareness exists only at $t^*$:

\begin{enumerate}
\item \textbf{Timing matters:} Interventions must reach people at their Moment 
of Truth, not before
\item \textbf{Context design:} The decision environment shapes which $A$ 
components are activated
\item \textbf{Repeated activation:} Single interventions fade; sustained change 
requires repeated $A$ activation at multiple $t^*$
\item \textbf{Habit formation:} The goal is often to move behavior from 
requiring $A$ (effortful) to not requiring $A$ (automatic)
\end{enumerate}

%------------------------------------------------------------------------------
\subsubsection{BEATRIX Module Reference}
%------------------------------------------------------------------------------

\begin{table}[htbp]
\centering
\caption{Section 11.2 BEATRIX References}
\label{tab:11-2-beatrix}
\renewcommand{\arraystretch}{1.3}
\begin{tabular}{@{}lll@{}}
\toprule
\textbf{Element} & \textbf{Module} & \textbf{Reference} \\
\midrule
Moment of Truth & BCM2\_05\_AWX\_7240\_02 & Sec\_Scope\_Kontext \\
Situativity Axiom & AWX-A2 & Fundamental Axiom \\
Context Integration & 8904\_KON & Context Module \\
\bottomrule
\end{tabular}
\end{table}

%%%%%%%%%%%%%%%%%%%%%%%%%%%%%%%%%%%%%%%%%%%%%%%%%%%%%%%%%%%%%%%%%%%%%%%%%%%%%%%
% END OF SECTION 11.2
%%%%%%%%%%%%%%%%%%%%%%%%%%%%%%%%%%%%%%%%%%%%%%%%%%%%%%%%%%%%%%%%%%%%%%%%%%%%%%%
